\section{Verification}
\label{sec:verification}

The construction and its parts were checked by direct execution, in the spirit of the machine checks that
accompany the tiling-specific automata. We describe the checks plainly. Each is part of the standing test
battery of the accompanying code \cite{pollard2026vibe} and passes.

\subsection{The locomotive}

A closed loop of track was built and the locomotive placed on it. Stepping the automaton, the head was recorded
at every beat. Over a full circuit the head visited every cell of the loop exactly once and returned to its
starting position, with the trailing cell always one behind it. This confirms Lemma~\ref{lem:motion} in
practice, the locomotive advances one cell per step and never reverses or stalls.

\subsection{The three switches}

Each switch was built with an entry lane, a trunk, and two branches, and the locomotive was driven through it.
A flip-flop switch entered actively from the trunk sent the locomotive to its selected branch and then flipped
its selection, and entered with the other selection it sent the locomotive to the other branch and flipped
back. A fixed switch sent the locomotive to its one branch and did not change. A memory switch entered passively
from a branch set its selection to that branch and then sent the locomotive to the trunk. In every case the
outcome matched the abstract railway switch of Section~\ref{sec:railway}, confirming Lemma~\ref{lem:switches}.

\subsection{The register machine}

At the level of the railway, registers and a sequencer were assembled into a two-register program, and the
single locomotive was run. A program for multiplication returned the product of its inputs, for several inputs,
in agreement with a reference register-machine interpreter run on the same programs. This confirms that correct
railway pieces compose into a working universal machine, which is the content of Proposition~\ref{prop:railway}
and the heart of Theorem~\ref{thm:main}.

\subsection{A tiling with no prior automaton}

To exhibit the uniformity concretely, the automaton was run on the octagrid, the tiling $\{8,3\}$ by regular
octagons three to a vertex. This is a hyperbolic regular tiling for which no universal cellular automaton had
been given. The cell-adjacency graph of the octagrid was generated, a closed track was laid along a cycle of
its cells, and the locomotive was set running. The locomotive traversed the octagrid track exactly as it
traverses a loop in any other tiling, because the rules are the same. The same automaton that runs the
pentagrid and the dodecagrid runs the octagrid, with no new rule and no new table.

\begin{corollary}[Octagrid]
\label{cor:octagrid}
The automaton of Theorem~\ref{thm:main} is weakly universal on the octagrid $\{8,3\}$, a tiling for which no
universal cellular automaton was previously known.
\end{corollary}

\subsection{Summary of the checks}

The locomotive moves correctly, the three switches behave as the railway demands, the register machine
computes, and the whole runs without change on a tiling outside the previously treated list. Together these are
the experimental content of Theorem~\ref{thm:main}.
